\section{So sánh Jetson Nano và Jetson Orin Nano}

So sánh giữa hai thiết bị giúp đánh giá tính di động của model và mức mở rộng
hiệu năng. Speedup được tính bằng tỷ số FPS Orin/Nano trên cùng cấu hình model;
latency phải dùng cùng định nghĩa và cùng pipeline.

\begin{table}[H]
\centering
\caption{So sánh hiệu năng giữa Jetson Nano và Jetson Orin Nano}
\label{tab:jetson-comparison}
\resizebox{\textwidth}{!}{%
\begin{tabular}{lccccc}
\toprule
\textbf{Model} & \textbf{Nano FPS} & \textbf{Orin FPS} & \textbf{Speedup} &
\textbf{Nano Latency} & \textbf{Orin Latency}\\
\midrule
FP32 & \cbs{} & \cbs{} & \cbs{} & \cbs{} & \cbs{}\\
FP16 & \cbs{} & \cbs{} & \cbs{} & \cbs{} & \cbs{}\\
INT8 & \cbs{} & \cbs{} & \cbs{} & \cbs{} & \cbs{}\\
Pruned + QAT & \cbs{} & \cbs{} & \cbs{} & \cbs{} & \cbs{}\\
\bottomrule
\end{tabular}}
\end{table}

Bảng~\ref{tab:jetson-comparison} cũng kiểm chứng việc Model Registry chọn đúng
artifact theo device class. Nếu cùng model version cần hai engine khác nhau, hai
artifact vẫn phải liên kết về cùng checkpoint và cấu hình tối ưu.